\scanpage{147}
% archive foliation: 146

\unsure{permet de se ramener}: un objet
\(P\) de la forme \(\underline h(X)\),
\deleted{\unsure{donc} \(u\) de la forme
\(\underline h(\alpha)\), où \(\alpha\in\operatorname{Fl}(X,X)\). Or comme}
\inserted{\(\operatorname{Fl}(P,P)\simeq\operatorname{Fl}(X,X)\), donc}
\(\operatorname{Fl}(P,P)\) est un \uline{anneau semi-simple} \(U\), \deleted{\unsure{Q}}
\deleted{Pour un tel anneau, tout \unsure{idéal opérant}
sur un objet d'une catégorie}
\inserted{Alors \(P\) et \(u\) se décomposent suivant les
composantes simples de \(U=\operatorname{Fl}(P,P)\) (i.e. de son
centre)}
\deleted{et on est ramené au cas où \(U\)}

Or lorsque un anneau de ce type opère
sur un objet d'une catégorie abélienne, alors
pour tout \(u\in U\), on peut trouver
\deleted{deux couples
d'idempotents orthogonaux complémentaires \(\pi,\pi'\) de \(U\),
\unsure{tels que}
\((\pi\pi'=\pi'\pi=0,\quad\pi+\pi'=1,\quad\pi^2=\pi,\quad\pi'^2=\pi')\)
tels que
et \(\varpi,\varpi'\) tels que}
\(u\pi'=0,\ \varpi'u=0,\)
\inserted{\(v\in U\) tel que
\(vu\) soit un \deleted{projecteur}\inserted{idempotent} \(\pi\)
\(uv\) soit \(\mathbin{-}\ \varpi\)}

\diagram{
\begin{tikzcd}[column sep=large,row sep=large]
P\simeq\pi P+\pi'P \arrow[r] & \pi P \arrow[d,"\unsure{\sim}"]\\
Q\simeq\varpi Q+\varpi'Q & \varpi Q \arrow[l]
\end{tikzcd}
}

\diagram{
\begin{tikzcd}[column sep=large,row sep=large]
P \arrow[r,"u"] \arrow[d] &
\unsure{Q} \arrow[l,bend right=40,"v"']\\
P' \arrow[r,"\simeq"] & Q' \arrow[u]
\end{tikzcd}
}

\[
\left\{
\begin{aligned}
vu&=\text{projecteur }\pi\\
uv&=\text{projecteur }\varpi\\
&\deleted{\unsure{vu}}
\end{aligned}
\right.
\]
\inserted{\figure{\begin{tikzpicture}[x=1cm,y=1cm]
\node at (0,0) {\(\begin{gathered}
\pi v=v,\quad u\pi=u\\
\varpi u=u,\quad v\varpi=v
\end{gathered}\)};
\draw (0,0) ellipse (2.25cm and .7cm);
\node[anchor=west] at (1.7,2.05) {et on a les relations};
\draw[->] (1.8,.5)..controls (2.5,1.3) and (4.9,.55)..(5.1,1.85);
\end{tikzpicture}}}
de cette façon,
gagne donc.

Pour les conditions de chaînes, on
utilise le fait que les \(\operatorname{Fl}(X,X)\) sont
de \deleted{\unsure{degr}} dim finie sur \(k\). .....

\diagram{
\begin{tikzcd}[column sep=large,row sep=large]
P'+P'' \arrow[d,shift left,"\unsure{\pi|\simeq}"'] & Q'+Q'' \arrow[d,shift left]\\
P' \arrow[u,shift left] \arrow[r,shift right] &
\unsure{Q'|Q} \arrow[u,shift left] \arrow[l,shift right]
\end{tikzcd}
}
